\documentclass[11pt,a4paper]{article}
\usepackage{fontspec,amsmath,amssymb,amsthm,microtype}
\setmainfont{lmroman10-regular.otf}[BoldFont=lmroman10-bold.otf,ItalicFont=lmroman10-italic.otf,BoldItalicFont=lmroman10-bolditalic.otf]
\setmonofont{lmmono10-regular.otf}[Scale=MatchLowercase]
\usepackage[margin=25mm]{geometry}
\usepackage[unicode,hidelinks]{hyperref}
\setlength{\emergencystretch}{3em}
\setlength{\parskip}{0.4em}
\setlength{\parindent}{0pt}
\allowdisplaybreaks[2]
\raggedbottom
\title{\texorpdfstring{The actual continuous Gibbs target and explicit exponential mixing}{The actual continuous Gibbs target and explicit exponential mixing}}
\author{LangevinLean mathematical exposition}
\date{}
\begin{document}\maketitle
\noindent Source module: \texttt{Kinetic\allowbreak{}Langevin/\allowbreak{}Continuous\allowbreak{}Gibbs.\allowbreak{}lean}. The accompanying HTML pairs this mathematical exposition with the exact Lean statements and records the scope of the formalization.


The actual canonical Brownian diffusion starts on an independent position–velocity–noise product. This module proves its exact Gibbs stationarity on that same source, rather than assuming a reference law. Combining it with the arbitrary-shift continuous coupling theorem gives unconditional transport to the true Gibbs target. The simple friction regime yields the explicit factor \(\sqrt{2}e^{-mt/\gamma}\).


\section{\texorpdfstring{Reassociate the genuine source and identify each thermal increment}{Reassociate the genuine source and identify each thermal increment}}


The measurable source map sends \((x,(v,B))\) to \(((v,x),B)\). Product associativity and exchange of independent position and velocity prove that this map preserves the exact phase-times-Brownian law. The continuous linear source pullbacks commute with the OU input integral, and the proved Brownian convolution formula identifies each actual filtered increment with the translated exponential kernel on the independent Brownian coordinate.



{\raggedright\noindent Lean declarations: \texttt{phase\allowbreak{}Source\allowbreak{}Projection}, \texttt{phase\allowbreak{}Source\allowbreak{}Projection\_\allowbreak{}preserving}, \texttt{ou\allowbreak{}Input\allowbreak{}Increment\_\allowbreak{}coe}.\par}

\section{\texorpdfstring{Prove the complete actual initial/\allowbreak{}noise law}{Prove the complete actual initial/noise law}}


Disjoint Brownian intervals give an actual iid stream of thermal Gaussian innovations. On the phase-times-Brownian product this stream is independent of the initial phase. Pulling that law back along the source reassociation and using the L\ensuremath{{}^{2}} representative identities proves the full initial-phase/\allowbreak{}entire-stream product law on the actual position–velocity–noise source.



{\raggedright\noindent Lean declarations: \texttt{initial\_\allowbreak{}ou\allowbreak{}Input\_\allowbreak{}joint\allowbreak{}Law}.\par}

\section{\texorpdfstring{Identify the exact stationary position law}{Identify the exact stationary position law}}


The selected continuous solution satisfies the genuine kinetic integral equation by construction. The actual initial/\allowbreak{}noise product law discharges every source premise of the proved splitting-limit stationarity theorem. This first proves the full actual phase has exactly the independent Gaussian-times-Gibbs target law. Projecting that phase law then gives the Gibbs position distribution at the terminal time. Strong convexity supplies the finite target second moment used in the construction.



{\raggedright\noindent Lean declarations: \texttt{phase\allowbreak{}Flow\_\allowbreak{}gibbs\_\allowbreak{}stationary}, \texttt{position\allowbreak{}Law\_\allowbreak{}gibbs}.\par}

\section{\texorpdfstring{Derive the full and simple explicit mixing bounds}{Derive the full and simple explicit mixing bounds}}


In the arbitrary-shift two-law theorem, start the reference position from the Gibbs distribution and use its newly proved stationarity. This produces the full explicit shifted coefficient with the true initial Wasserstein distance. For shift zero and \(m+M\le\gamma^2\), the maximum in the exponent equals \(-m\); the prefactor simplifies to \(\sqrt2\). Finite L\ensuremath{{}^{2}} output and Gibbs moments also prove the actual transport distance is finite, ensuring the real-valued endpoint never conceals an infinite distance.



{\raggedright\noindent Lean declarations: \texttt{wasserstein2\_\allowbreak{}gibbs\_\allowbreak{}general}, \texttt{wasserstein2\_\allowbreak{}gibbs}, \texttt{wasserstein2\_\allowbreak{}gibbs\_\allowbreak{}ne\_\allowbreak{}top}.\par}

\end{document}
